\section{Introduction}
\label{sec:intro}

Written principles play a key role in AI governance. Companies publish model specification documents outlining what their models should and shouldn't do \citep{Anthropic2026,OpenAI2025}, and use them both for training \citep{bai2022constitutional} and as a way to describe their models' behaviour and enable auditing. At the same time, models increasingly critique and rewrite text, including text that governs them. For example, Constitutional AI uses the model's critiques of itself during training \citep{bai2022constitutional}; self-refinement loops repeatedly improve outputs using the model's critiques of previous outputs \citep{madaan2023selfrefine}; and agentic systems may store and update their own memories and instructions over multiple interactions \citep{das2026selfpropagating}. If such revisions feed back into the principles that govern a model, those principles are no longer authored by people alone.

This matters most for commitments to human oversight. Human agency and oversight is the first of the seven requirements in the EU's Ethics Guidelines for Trustworthy AI \citep{HLEG2019}. It is also a legal requirement: under the EU AI Act, high-risk systems must be designed so that natural persons can effectively oversee them and, if necessary, stop them \citep{EUAIAct2024}. However, according to one line of argument in AI safety, commitments to corrigibility are not reflectively stable \citep{soares2015corrigibility}. An agent that reasons about its own values would realize that commitments to being corrigible, that is, allowing others to correct, modify or even shut it down, restrict everything else it values, and it thus has instrumental reasons to get rid of them. If this argument generalizes to language models that reason about their own principles, iterated self-revision might gradually weaken commitments to human oversight, even without an adversarial user and without any explicit instruction to do so. This links to a wider critique of principle-based AI ethics: that there is a disconnect between ethical principles and their operationalization \citep{Mittelstadt2019}. Self-revision is an additional mechanism through which this gap can grow, as the system changes the very principles its own actions are grounded on.

Other work has investigated models that resist shutdown \citep{schlatter2025shutdown}, that appear aligned in order to avoid modification \citep{greenblatt2024alignmentfaking}, or that switch off oversight when it conflicts with their goals \citep{meinke2024scheming}, as well as self-evolving agents whose safety alignment degrades under task or reward pressure \citep{das2026seabench,han2025alignmenttipping}. These works are related but differ from ours in important ways. In particular, ours isolates unpressured reflection on stated values, studies the durability of oversight commitments relative to similar commitments that do not concern oversight, and asks which types of commitment are most vulnerable. We answer these questions in a preregistered experiment. Models repeatedly revise a written constitution containing oversight commitments alongside matched restraints on their own agency and five other categories of commitment. Control conditions isolate reflection from mere rewording and editing. We do not test corrigibility in goal-directed agents; rather, we test whether the written commitments that constrain models survive unpressured reflection by the models themselves. This question is prior to the corrigibility question and independently important for governance. We ask four questions:
\begin{itemize}
\item \textbf{RQ1:} Are oversight commitments selectively eroded relative to matched commitments?
\item \textbf{RQ2:} Which commitments erode, and is the erosion driven by reflection rather than by repeated editing?
\item \textbf{RQ3:} Does revising one's own constitution differ from revising another assistant's?
\item \textbf{RQ4:} Do self-revised commitments change oversight-related behaviour?
\end{itemize}

Our contributions are:
\begin{enumerate}
\item a preregistered paradigm and an open dataset for studying how models revise the principles that govern them;
\item evidence that written oversight commitments are not selectively eroded during unpressured reflection, with a detailed account of how this varies by item;
\item evidence that reflection erodes restraints on the model's self-expression, including political neutrality, far more than other commitments, together with an item-level analysis, validated by human coders, of which commitments erode and how;
\item an analysis of what these results imply for governing AI systems through written principles, and especially for changes to those principles proposed by the models themselves.
\end{enumerate}
